% ==========================================
% Glossario C6 -- Casi d'uso
% ==========================================
% Documento del gruppo Bit Crusher, compilabile con pdfLaTeX, autonomo
% come il modello del Glossario: non include altri sorgenti.
% Le voci sono scritte con \VoceGlossario{id}{nome}{definizione}, come nel
% Glossario di progetto, e la pagina glossario.html le legge da questo file.
% Livelli, ereditarieta' e regole: vedi l'Introduzione e
% docs/requisiti/glossari.yml.

\documentclass[a4paper, 11pt]{article}

% --- PACCHETTI ---
\usepackage[utf8]{inputenc}
\usepackage[T1]{fontenc}
\usepackage[italian]{babel}
\usepackage[margin=2.5cm]{geometry}
\usepackage{graphicx}
\usepackage{tabularx}
\usepackage{array}
\usepackage{fancyhdr}
\usepackage[hidelinks, hypertexnames=false]{hyperref}

% --- METADATI ---
\newcommand{\VersioneDocumento}{0.1.0}
\newcommand{\DataDocumento}{2026-11-26}
\newcommand{\RedattoriDocumento}{Alex Singh}
\newcommand{\VerificatoriDocumento}{[Nome Cognome]}
\newcommand{\ApprovatoreDocumento}{[Nome Cognome]}

\hypersetup{
    pdftitle={Glossario C6 -- Casi d'uso},
    pdfauthor={Bit Crusher},
    pdfsubject={Glossario del capitolato C6 per livello}
}

% --- HEADER E FOOTER ---
\pagestyle{fancy}
\fancyhf{}
\fancyhead[L]{Bit Crusher}
\fancyhead[R]{Glossario C6, casi d'uso}
\fancyfoot[C]{\thepage}
\setlength{\headheight}{14pt}
\renewcommand{\headrulewidth}{0.4pt}

% ==========================================
% COMANDI PER LE VOCI E I COLLEGAMENTI (come nel Glossario di progetto)
% ==========================================
% \VoceGlossario{id}{nome}{definizione}: stampa la voce e crea l'ancora gls-<id>.
% Una chiave ripetuta produce un errore di compilazione.
\newcommand{\VoceGlossario}[3]{%
    \ifcsname bitcrusherglossario@#1\endcsname
        \PackageError{Glossario}{Identificatore duplicato: #1}%
            {Assegnare un identificatore diverso a ogni voce.}%
    \else
        \expandafter\gdef\csname bitcrusherglossario@#1\endcsname{}%
        \par\addvspace{\baselineskip}%
        \noindent\hypertarget{gls-#1}{\textbf{#2}}\par
        \nopagebreak\smallskip
        \noindent #3\par
    \fi
}
% \RifGlossario{id}{testo}: rimando a una voce di QUESTO documento.
\newcommand{\RifGlossario}[2]{%
    \hyperlink{gls-#1}{#2}%
}
% \RifGlossarioWeb{URL del PDF}{id}{testo}: rimando a una voce di un ALTRO
% glossario pubblicato (Glossario di progetto o livello C6 precedente).
\newcommand{\RifGlossarioWeb}[3]{%
    \href{#1\#nameddest=gls-#2}{#3}%
}

\begin{document}

% ==========================================
% FRONTESPIZIO
% ==========================================
\begin{titlepage}
    \centering
    \begin{minipage}{0.45\textwidth}
        \centering
        \IfFileExists{UnipdLogoRosso.png}%
            {\includegraphics[width=0.7\textwidth]{UnipdLogoRosso.png}}%
            {\fbox{\parbox[c][2.5cm][c]{0.7\textwidth}{\centering Logo UniPD}}}
    \end{minipage}
    \hfill
    \begin{minipage}{0.45\textwidth}
        \centering
        \IfFileExists{BitCrusherTeamNero.png}%
            {\includegraphics[width=0.7\textwidth]{BitCrusherTeamNero.png}}%
            {\fbox{\parbox[c][2.5cm][c]{0.7\textwidth}{\centering Logo Bit Crusher}}}
    \end{minipage}

    \vspace{1.5cm}
    {\Large\textbf{UNIVERSITÀ DEGLI STUDI DI PADOVA}\par}
    \vspace{0.3cm}
    {\large Corso di Laurea in Informatica\par}
    \vspace{0.3cm}
    {\large Insegnamento di Ingegneria del Software\par}
    \vspace{0.3cm}
    {\large Anno Accademico 2026/2027\par}

    \vspace{1.5cm}
    {\Huge\textbf{Glossario C6 -- Casi d'uso}\par}
    \vspace{0.5cm}
    {\Large Versione \VersioneDocumento\par}

    \vspace{1.5cm}
    \begin{flushleft}
        \large
        \textbf{Gruppo:} Bit Crusher \\
        \vspace{0.1cm}
        \textbf{Email:} \href{mailto:bitcrusherteam@gmail.com}{bitcrusherteam@gmail.com} \\
        \vspace{0.1cm}
        \textbf{Redattori:} \RedattoriDocumento \\
        \vspace{0.1cm}
        \textbf{Verificatori:} \VerificatoriDocumento \\
        \vspace{0.1cm}
        \textbf{Data ultima modifica:} \DataDocumento
    \end{flushleft}
    \vfill
\end{titlepage}

% ==========================================
% REGISTRO DELLE MODIFICHE (fuori dall'indice, dal piu' recente)
% ==========================================
\section*{Registro delle modifiche}
\begin{center}
    \renewcommand{\arraystretch}{1.5}
    \begin{tabularx}{\textwidth}{|l|l|>{\raggedright\arraybackslash}X|>{\raggedright\arraybackslash}X|>{\raggedright\arraybackslash}X|}
        \hline
        \textbf{Versione} & \textbf{Data} & \textbf{Autore} & \textbf{Verificatore} & \textbf{Descrizione} \\ \hline
        0.1.0 & 2026-11-26 & Alex Singh & [Nome Cognome] & Prima stesura delle voci. \\ \hline
    \end{tabularx}
\end{center}
\newpage

\tableofcontents
\newpage

\section{Introduzione}

\subsection{Scopo del documento}
Il presente glossario definisce i termini del dominio del capitolato C6 usati nei casi d'uso dell'Analisi dei Requisiti: gli attori, gli oggetti e le operazioni che il prodotto offre, descritti dall'esterno e senza riferimento al modo in cui saranno realizzati. Raccoglie anche i termini dei bisogni ricavati dalle fonti, quando servono a darne un significato univoco.

\subsection{Livelli ed ereditarietà}
I glossari del capitolato C6 sono tre, uno per livello: casi d'uso, requisiti e progettazione. Ogni livello eredita le voci dei livelli precedenti e aggiunge le proprie. Un testo di un livello usa le voci del proprio livello e di quelli precedenti, mai quelle dei successivi: un caso d'uso non usa termini dei requisiti, un requisito non usa termini di progettazione.
\begin{center}
Glossario di progetto $\rightarrow$ C6 casi d'uso $\rightarrow$ C6 requisiti $\rightarrow$ C6 progettazione
\end{center}
Questo glossario eredita le voci del \href{https://bitcrusher-team.github.io/Progetto_SWE_26-27/docs/compiled/glossario/Glossario.pdf}{Glossario di progetto}, che non ripete. Le sue voci sono a loro volta ereditate dal glossario dei requisiti e da quello di progettazione.

\subsection{Regole di univocità}
\begin{itemize}
    \item Ogni voce ha un identificatore unico in tutti i glossari: un livello non ridefinisce una voce ereditata.
    \item In una stessa catena di livelli due voci non hanno lo stesso nome.
    \item Una definizione usa solo voci del proprio livello o dei precedenti, oppure parole del linguaggio comune nel loro significato comune.
    \item Le voci del Glossario di progetto che descrivono una soluzione, come \textit{Embedding}, \textit{TTL} o \textit{Soglia di similarità}, appartengono alla progettazione: i casi d'uso e i requisiti non le usano.
    \item Lo script \texttt{strumenti/verifica\_requisiti.py} controlla queste regole sui glossari e sugli elementi Doorstop.
\end{itemize}

\subsection{Convenzioni di lettura}
\begin{itemize}
    \item \textbf{Ordinamento:} le voci sono in ordine alfabetico, raggruppate per lettera iniziale.
    \item \textbf{Rimandi:} i riferimenti ad altre voci sono in corsivo e cliccabili. Un rimando a una voce di questo glossario porta alla voce; un rimando a un altro glossario apre il suo PDF pubblicato sulla voce.
\end{itemize}

\newpage
% ==========================================
% TERMINI: generati in ordine alfabetico; una voce per concetto.
% Rimandi: \RifGlossario per questo documento, \RifGlossarioWeb per gli altri glossari.
% ==========================================

\section{A}

\VoceGlossario{accesso-differenziato}{Accesso differenziato}{
    Proprietà di una \RifGlossario{richiesta-di-primo-turno}{\textit{richiesta di primo turno}} la cui risposta corretta dipende da dati che non tutti gli utenti del suo \RifGlossario{ambito-di-condivisione}{\textit{ambito di condivisione}} possono consultare: posta da utenti diversi, la stessa richiesta può avere risposte corrette diverse. Esempio: una domanda a cui l'agente risponde consultando documenti riservati del reparto HR.
}

\VoceGlossario{ambito-di-condivisione}{Ambito di condivisione}{
    Insieme degli utenti a cui può essere restituito il contenuto di un \RifGlossario{elemento-conservato}{\textit{elemento conservato}}. Ha uno di tre valori: \textit{ambito utente}, il solo utente d'origine; \textit{ambito team}, gli utenti del team d'origine; \textit{ambito globale}, tutti gli utenti (vedi \RifGlossario{richiesta-d-origine}{\textit{Richiesta d'origine}}).
}

\VoceGlossario{amministratore-litellm}{Amministratore LiteLLM}{
    Persona con il ruolo di amministratore nell'istanza di \RifGlossarioWeb{https://bitcrusher-team.github.io/Progetto_SWE_26-27/docs/compiled/glossario/Glossario.pdf}{litellm}{\textit{LiteLLM}} di Var Group. Ne monitora costi e consumi e, nel prodotto, gestisce le impostazioni di riuso e di conservazione, la \RifGlossario{retention}{\textit{retention}} e gli \RifGlossario{elemento-conservato}{\textit{elementi conservati}}. Da non confondere con l'\RifGlossarioWeb{https://bitcrusher-team.github.io/Progetto_SWE_26-27/docs/compiled/glossario/Glossario.pdf}{amministratore}{\textit{Amministratore}}, ruolo del progetto didattico.
}

\section{C}

\VoceGlossario{client-claude}{Client Claude}{
    Applicazione Claude di Anthropic, nella versione desktop, usata dagli utenti di Var Group. Invia le \RifGlossario{richiesta-di-generazione}{\textit{richieste di generazione}} dei suoi utenti a LiteLLM, e quindi al sistema, di cui è l'attore che chiede le risposte.
}

\VoceGlossario{confidenza-di-sensibilita}{Confidenza di sensibilità}{
    Numero compreso tra 0 e 1 che il sistema associa a ogni \RifGlossario{elemento-conservato}{\textit{elemento conservato}} per stimare la probabilità che la richiesta d'origine o la risposta contengano \RifGlossario{dato-sensibile}{\textit{dati sensibili}}.
}

\VoceGlossario{conservazione}{Conservazione}{
    Operazione con cui il sistema registra la risposta a una \RifGlossario{richiesta-di-primo-turno}{\textit{richiesta di primo turno}} come \RifGlossario{elemento-conservato}{\textit{elemento conservato}}.
}

\VoceGlossario{conversazione}{Conversazione}{
    Sequenza di messaggi scambiati tra un utente e il modello nel \RifGlossario{client-claude}{\textit{client Claude}}. Ogni \RifGlossario{richiesta-di-generazione}{\textit{richiesta di generazione}} di una conversazione contiene tutti i messaggi precedenti.
}

\VoceGlossario{costo-evitato}{Costo evitato}{
    Per una richiesta servita con \RifGlossario{riuso}{\textit{riuso}}, costo che la richiesta avrebbe comportato se fosse stata inoltrata al \RifGlossarioWeb{https://bitcrusher-team.github.io/Progetto_SWE_26-27/docs/compiled/glossario/Glossario.pdf}{provider}{\textit{provider}}. Si calcola come il \RifGlossario{costo-sostenuto}{\textit{costo sostenuto}} per generare la risposta dell'elemento riusato.
}

\VoceGlossario{costo-sostenuto}{Costo sostenuto}{
    Costo addebitato a Var Group dal \RifGlossarioWeb{https://bitcrusher-team.github.io/Progetto_SWE_26-27/docs/compiled/glossario/Glossario.pdf}{provider}{\textit{provider}} per una \RifGlossario{richiesta-di-generazione}{\textit{richiesta di generazione}} inoltrata, calcolato sui \RifGlossarioWeb{https://bitcrusher-team.github.io/Progetto_SWE_26-27/docs/compiled/glossario/Glossario.pdf}{token}{\textit{token}} consumati con le tariffe applicate a Var Group.
}

\section{D}

\VoceGlossario{dato-sensibile}{Dato sensibile}{
    Nel prodotto: qualsiasi dato personale nel senso del \RifGlossarioWeb{https://bitcrusher-team.github.io/Progetto_SWE_26-27/docs/compiled/glossario/Glossario.pdf}{gdpr}{\textit{GDPR}} o informazione aziendale riservata di Var Group. Il termine è più ampio delle «categorie particolari di dati personali» dell'articolo 9 del GDPR, con cui non va confuso.
}

\section{E}

\VoceGlossario{elemento-conservato}{Elemento conservato}{
    Risposta che il sistema ha registrato per poterla restituire ad altre \RifGlossario{richiesta-di-primo-turno}{\textit{richieste di primo turno}}, insieme alla \RifGlossario{richiesta-d-origine}{\textit{richiesta d'origine}}, all'utente e al team d'origine, al modello che l'ha generata, all'\RifGlossario{ambito-di-condivisione}{\textit{ambito di condivisione}}, all'istante di conservazione e alla \RifGlossario{scadenza}{\textit{scadenza}}. Negli appunti di Var Group è l'elemento in cache.
}

\VoceGlossario{elemento-valido}{Elemento valido}{
    \RifGlossario{elemento-conservato}{\textit{Elemento conservato}} la cui \RifGlossario{scadenza}{\textit{scadenza}} non è ancora trascorsa e che non è stato eliminato.
}

\section{I}

\VoceGlossario{impostazione-di-conservazione}{Impostazione di conservazione}{
    Scelta di un \RifGlossario{utente}{\textit{utente}}, attiva o disattiva, che stabilisce se le risposte alle sue \RifGlossario{richiesta-di-primo-turno}{\textit{richieste di primo turno}} possono diventare \RifGlossario{elemento-conservato}{\textit{elementi conservati}}. Il sistema applica l'\RifGlossario{impostazione-effettiva}{\textit{impostazione effettiva}}.
}

\VoceGlossario{impostazione-di-riuso}{Impostazione di riuso}{
    Scelta di un \RifGlossario{utente}{\textit{utente}}, attiva o disattiva, che stabilisce se alle sue \RifGlossario{richiesta-di-primo-turno}{\textit{richieste di primo turno}} può essere restituito il contenuto di un \RifGlossario{elemento-conservato}{\textit{elemento conservato}}. Il sistema applica l'\RifGlossario{impostazione-effettiva}{\textit{impostazione effettiva}}.
}

\VoceGlossario{impostazione-effettiva}{Impostazione effettiva}{
    Valore che il sistema applica a un \RifGlossario{utente}{\textit{utente}} per il riuso e, separatamente, per la conservazione. È attiva se e solo se l'impostazione dell'utente è attiva e l'\RifGlossario{amministratore-litellm}{\textit{amministratore LiteLLM}} non ha disattivato quella funzione né per l'utente né per il suo \RifGlossario{team}{\textit{team}}. Ne segue che l'amministratore può disattivare una funzione per un utente, ma non attivarla al suo posto.
}

\section{L}

\VoceGlossario{latenza-aggiunta}{Latenza aggiunta}{
    Tempo che il sistema aggiunge a una \RifGlossario{richiesta-di-generazione}{\textit{richiesta di generazione}}. Per una richiesta inoltrata al \RifGlossarioWeb{https://bitcrusher-team.github.io/Progetto_SWE_26-27/docs/compiled/glossario/Glossario.pdf}{provider}{\textit{provider}}: tempo trascorso nel sistema tra la ricezione della richiesta e il suo inoltro, più il tempo tra l'arrivo del primo frammento della risposta del provider e la sua restituzione al \RifGlossario{client-claude}{\textit{client Claude}}. Per una richiesta servita con \RifGlossario{riuso}{\textit{riuso}}: tempo tra la ricezione della richiesta e la restituzione del primo frammento della risposta.
}

\VoceGlossario{lettore-in-ombra}{Lettore in ombra}{
    Modalità di funzionamento, attivabile per \RifGlossario{team}{\textit{team}}, in cui il sistema stabilisce per ogni \RifGlossario{richiesta-di-primo-turno}{\textit{richiesta di primo turno}} quale \RifGlossario{elemento-conservato}{\textit{elemento conservato}} riuserebbe, ma inoltra comunque la richiesta al \RifGlossarioWeb{https://bitcrusher-team.github.io/Progetto_SWE_26-27/docs/compiled/glossario/Glossario.pdf}{provider}{\textit{provider}} e restituisce la risposta generata. Registra la richiesta, l'elemento individuato e la risposta generata, in modo da poterli confrontare. In inglese si dice \textit{shadow mode}.
}

\section{M}

\VoceGlossario{mcp}{MCP (Model Context Protocol)}{
    Protocollo aperto con cui un'applicazione, come il \RifGlossario{client-claude}{\textit{client Claude}}, si collega a server esterni che le offrono \RifGlossario{strumento}{\textit{strumenti}} e dati.
}

\section{O}

\VoceGlossario{one-shot}{One-shot}{
    Termine con cui Var Group indica le \RifGlossario{richiesta-di-primo-turno}{\textit{richieste di primo turno}}. Da non confondere con il \textit{one-shot prompting}, in cui il prompt contiene un solo esempio del compito. Vedi anche \RifGlossarioWeb{https://bitcrusher-team.github.io/Progetto_SWE_26-27/docs/compiled/glossario/Glossario.pdf}{single-shot}{\textit{Single-shot}}.
}

\section{R}

\VoceGlossario{ragionamento}{Ragionamento (reasoning)}{
    Parte dell'elaborazione di un modello che precede la risposta. Nei modelli Claude è restituita in blocchi separati dal testo della risposta, ciascuno legato a tutto ciò che lo precede nella richiesta, ed è fatturata come \RifGlossarioWeb{https://bitcrusher-team.github.io/Progetto_SWE_26-27/docs/compiled/glossario/Glossario.pdf}{token}{\textit{token}} in uscita.
}

\VoceGlossario{retention}{Retention}{
    Durata per cui un nuovo \RifGlossario{elemento-conservato}{\textit{elemento conservato}} resta riusabile: la sua \RifGlossario{scadenza}{\textit{scadenza}} è l'istante di conservazione più la retention in vigore. È il termine usato da Var Group; la sua realizzazione, ad esempio con un \RifGlossarioWeb{https://bitcrusher-team.github.io/Progetto_SWE_26-27/docs/compiled/glossario/Glossario.pdf}{ttl}{\textit{TTL}}, è materia di progettazione.
}

\VoceGlossario{richiesta-d-origine}{Richiesta d'origine}{
    \RifGlossario{richiesta-di-primo-turno}{\textit{Richiesta di primo turno}} la cui risposta è diventata un \RifGlossario{elemento-conservato}{\textit{elemento conservato}}. L'utente e il team che l'hanno inviata sono l'\textit{utente d'origine} e il \textit{team d'origine} dell'elemento.
}

\VoceGlossario{richiesta-di-generazione}{Richiesta di generazione}{
    Richiesta inviata dal \RifGlossario{client-claude}{\textit{client Claude}} per ottenere una risposta da un modello. Contiene il modello scelto, le istruzioni di sistema, gli \RifGlossario{strumento}{\textit{strumenti}} disponibili e tutti i messaggi della \RifGlossario{conversazione}{\textit{conversazione}} fino a quel momento.
}

\VoceGlossario{richiesta-di-primo-turno}{Richiesta di primo turno}{
    \RifGlossario{richiesta-di-generazione}{\textit{Richiesta di generazione}} che contiene un solo messaggio dell'utente e nessun messaggio del modello. Le richieste che proseguono una risposta dopo l'uso di uno \RifGlossario{strumento}{\textit{strumento}} contengono messaggi del modello, quindi non sono di primo turno. Var Group le chiama \RifGlossario{one-shot}{\textit{one-shot}}.
}

\VoceGlossario{richiesta-non-permanente}{Richiesta non permanente}{
    \RifGlossario{richiesta-di-primo-turno}{\textit{Richiesta di primo turno}} la cui risposta corretta può cambiare entro un giorno, perché dipende dall'istante in cui è posta. Esempio: «che tempo fa oggi a Padova?».
}

\VoceGlossario{richieste-equivalenti}{Richieste equivalenti}{
    Due \RifGlossario{richiesta-di-primo-turno}{\textit{richieste di primo turno}} rivolte allo stesso modello sono equivalenti quando ogni risposta corretta all'una è corretta anche per l'altra, e viceversa. L'equivalenza riguarda il significato, non il testo: richieste scritte in modo diverso possono essere equivalenti. Stabilirla spetta a valutatori umani; il giudizio del sistema è il \RifGlossario{riconoscimento-di-equivalenza}{\textit{riconoscimento di equivalenza}}.
}

\VoceGlossario{riconoscimento-di-equivalenza}{Riconoscimento di equivalenza}{
    Giudizio con cui il sistema stabilisce che una \RifGlossario{richiesta-di-primo-turno}{\textit{richiesta di primo turno}} è \RifGlossario{richieste-equivalenti}{\textit{equivalente}} alla \RifGlossario{richiesta-d-origine}{\textit{richiesta d'origine}} di un \RifGlossario{elemento-conservato}{\textit{elemento conservato}}. Può sbagliare: un riconoscimento errato seguito da riuso è un \RifGlossarioWeb{https://bitcrusher-team.github.io/Progetto_SWE_26-27/docs/compiled/glossario/Glossario.pdf}{falso-hit}{\textit{falso hit}}. La sua prudenza dipende dalla \RifGlossario{severita-del-riconoscimento}{\textit{severità del riconoscimento}}.
}

\VoceGlossario{risparmio-netto}{Risparmio netto}{
    Per un intervallo di tempo: somma dei \RifGlossario{costo-evitato}{\textit{costi evitati}} meno i costi che il sistema sostiene nello stesso intervallo per il \RifGlossario{riconoscimento-di-equivalenza}{\textit{riconoscimento di equivalenza}} e per la \RifGlossario{conservazione}{\textit{conservazione}}.
}

\VoceGlossario{riuso}{Riuso}{
    Restituzione al \RifGlossario{client-claude}{\textit{client Claude}}, come risposta a una \RifGlossario{richiesta-di-primo-turno}{\textit{richiesta di primo turno}}, del contenuto di un \RifGlossario{elemento-conservato}{\textit{elemento conservato}} al posto di una risposta generata dal modello per quella richiesta. Una richiesta servita con riuso è un \RifGlossarioWeb{https://bitcrusher-team.github.io/Progetto_SWE_26-27/docs/compiled/glossario/Glossario.pdf}{cache-hit}{\textit{cache hit}}.
}

\section{S}

\VoceGlossario{scadenza}{Scadenza}{
    Istante dopo il quale il contenuto di un \RifGlossario{elemento-conservato}{\textit{elemento conservato}} non può più essere restituito.
}

\VoceGlossario{segnalazione-di-riuso}{Segnalazione di riuso}{
    Testo che il sistema aggiunge a una risposta restituita con \RifGlossario{riuso}{\textit{riuso}} per indicare all'utente che proviene da un \RifGlossario{elemento-conservato}{\textit{elemento conservato}}.
}

\VoceGlossario{severita-del-riconoscimento}{Severità del riconoscimento}{
    Parametro, impostato dall'\RifGlossario{amministratore-litellm}{\textit{amministratore LiteLLM}}, che stabilisce quanto deve essere prudente il \RifGlossario{riconoscimento-di-equivalenza}{\textit{riconoscimento di equivalenza}}: a una severità maggiore corrispondono meno riusi e meno \RifGlossarioWeb{https://bitcrusher-team.github.io/Progetto_SWE_26-27/docs/compiled/glossario/Glossario.pdf}{falso-hit}{\textit{falsi hit}}. Il capitolato la chiama \RifGlossarioWeb{https://bitcrusher-team.github.io/Progetto_SWE_26-27/docs/compiled/glossario/Glossario.pdf}{soglia-di-similarita}{\textit{soglia di similarità}}, nome che descrive una possibile realizzazione.
}

\VoceGlossario{skill}{Skill}{
    Pacchetto di istruzioni, a volte con script e risorse, che il \RifGlossario{client-claude}{\textit{client Claude}} mette a disposizione del modello e che il modello carica nel proprio contesto quando lo ritiene pertinente alla richiesta.
}

\VoceGlossario{strumento}{Strumento (tool)}{
    Funzione esterna che il modello può chiedere al \RifGlossario{client-claude}{\textit{client Claude}} di eseguire durante una risposta, ad esempio una ricerca in un archivio. Il risultato torna al modello in una \RifGlossario{richiesta-di-generazione}{\textit{richiesta di generazione}} successiva. Gli strumenti possono essere offerti da server \RifGlossario{mcp}{\textit{MCP}}.
}

\section{T}

\VoceGlossario{team}{Team}{
    Gruppo di utenti definito in \RifGlossarioWeb{https://bitcrusher-team.github.io/Progetto_SWE_26-27/docs/compiled/glossario/Glossario.pdf}{litellm}{\textit{LiteLLM}}, a cui si attribuiscono costi e impostazioni. Nel prodotto si assume che ogni utente appartenga a uno e un solo team; l'ipotesi va confermata con Var Group.
}

\section{U}

\VoceGlossario{utente}{Utente}{
    Dipendente di Var Group che usa il \RifGlossario{client-claude}{\textit{client Claude}} ed è registrato in \RifGlossarioWeb{https://bitcrusher-team.github.io/Progetto_SWE_26-27/docs/compiled/glossario/Glossario.pdf}{litellm}{\textit{LiteLLM}}.
}

\end{document}
